%----------------------------------------------------------------------------------------------------
% START OF CHAPTER
%----------------------------------------------------------------------------------------------------
\chapter{May 4}
\paragraph{Topics}
\textit{Direct and indirect object pronouns, Sentence review}

\section{Grammar} \label{0504-Grammar}
\subsection{Direct object pronouns}

A direct object pronoun replaces the noun object in a sentence. `I eat gelato' $\rightarrow$ `I eat it'.
They stand in for whatever thing the verb is acting on.
\begin{mdframed}
	Object pronouns in Italian replace nouns to avoid repetition. They answer the questions 
	`what?' or `whom?' in relation to the verb.
\end{mdframed}

\begin{table}[ht]
	\centering
	\begin{tabular}{ r | l }
	                         & pronoun   \\
		\hline
		\textbf{me}          & \iti{mi}  \\
		\textbf{you}         & \iti{ti}  \\
		\textbf{him/it}      & \iti{lo}  \\
		\textbf{her/it}      & \iti{la}  \\
		\textbf{us}          & \iti{ci}  \\
		\textbf{you pl.}     & \iti{vi}  \\
		\textbf{them, masc.} & \iti{li}  \\
		\textbf{them, fem.}  & \iti{le}
	\end{tabular}
	\caption{Direct object pronouns}
	\label{tab:tDirectObjectPronouns}
\end{table}
	
\paragraph{Examples}
\begin{itemize}
	\item \iti{Mangio gelato.} $\rightarrow$ \iti{Lo mangio.}
	\item \iti{Mangio pizza.} $\rightarrow$ \iti{La mangia.}
\end{itemize}
	
In the past (perfect) tense, the past participle must be modified to match the gender/plurality
of the pronoun, as if it is working as an adjective.
	
\paragraph{Examples}
\begin{itemize}
	\item{I bought the car. I bought it} \\
	\iti{Ho comprato la macchina. L'ho comprata.}
		
	\item{Ho visto che tu l'hai comprato già} \\
	\iti{I saw that you had already bought it. (masc.)}
		
	\item {Ho visto che tu l'hai comprata già} \\
	\iti{I saw that you had already bought it. (fem.)}
\end{itemize}
	
\subsection{Indirect object pronouns}
These pronouns replace the recipient of the action of the verb, e.g. `I told \textbf{you}'.
	
\begin{mdframed}
	Indirect object pronouns replace the noun that is the recipient of the action, 
	answering `to whom?' or `for whom?' the action is done.
\end{mdframed}
	
\begin{table}[ht]
	\centering
	\begin{tabular}{ r l | l }
	               & English           & Italian    \\
		\hline
			to/for & \textbf{me}       & \iti{mi}   \\
			       & \textbf{you}      & \iti{ti}   \\
                   & \textbf{him}      & \iti{gli}  \\
                   & \textbf{her}      & \iti{le}   \\
                   & \textbf{us}       & \iti{ci}   \\
                   & \textbf{you pl.}  & \iti{vi}   \\
                   & \textbf{them}     & \iti{gli}
	\end{tabular}
	\caption{Indirect object pronouns}
	\label{tab:tIndirectObjectPronouns}
\end{table}
	
Tip: if in doubt, decide what the verb is acting upon (the direct object), and 
to whom or for whom the action is done (the indirect object).
	
\paragraph{Examples}
\begin{itemize}
	\item{I told you that I bought the ice cream for tonight.} \\
	\iti{Ti ho detto che ho comprato il gelato per questa sera.}
		
	\item{I told you already that I worked a lot last night.} \\
	\iti{Ti ho detto già che ho lavorato molto ieri sera.}
\end{itemize}
	
\section{Sentence review} \label{0504-Sentences}
\begin{itemize}
	\item{I'm happy because my wife already bought a table for our living room. We needed this table
		for when we have guests for dinner. The table is big and is for twelve people.
		I'm happy that she has already bought it.} \\
	\iti{Sono felice perché mia moglie ha già comprato un tavolo per il nostro soggiorno. 
		Abbiamo bisogno di questo tavolo quando abbiamo ospiti per cena. Il tavolo è grande 
		ed è per dodici persone. Sono felice che lo ha comprato già.}
	
	\item{Hi Giorgio! I called you last night but you didn't answer. I hope everything is OK.
		I'm working a lot these weeks. I told you that I wanted to go on vacation with you and your family.
		Let's talk about that soon. See you soon!} \\
	\iti{Ciao Giorgio! Ho ti chiamato ieri sera, ma tu non hai risposto. Spero tutto bene. Sto lavorando
		molto queste settimane. Ti ho detto che voglio andare in vacanza con te e la tua famiglia.
		Parliamo di questo presto. A presto!}
\end{itemize}

%----------------------------------------------------------------------------------------------------
% END OF CHAPTER
%----------------------------------------------------------------------------------------------------